\begin{tikzpicture}[font=\small,line cap=round,line join=round]\draw[geoTrait,fill=otNavyTint](0.0,0)--++(.9,0)--++(-.3,.9)--++(-.3,-.35)--cycle;\node[below,font=\scriptsize]at(0.45,-.05){1};\draw[geoTrait,fill=otNavyTint](1.6,0)--++(.9,0)--++(-.3,.9)--++(-.3,-.35)--cycle;\node[below,font=\scriptsize]at(2.0500000000000003,-.05){2};\draw[geoTrait,fill=otNavyTint](3.2,0)--++(.9,0)--++(-.3,.9)--++(-.3,-.35)--cycle;\node[below,font=\scriptsize]at(3.6500000000000004,-.05){3};\draw[geoTrait,fill=otNavyTint](4.800000000000001,0)--++(.9,0)--++(-.3,.9)--++(-.3,-.35)--cycle;\node[below,font=\scriptsize]at(5.250000000000001,-.05){4};\draw[geoTrait,fill=otNavyTint](6.4,0)--++(.9,0)--++(-.3,.9)--++(-.3,-.35)--cycle;\node[below,font=\scriptsize]at(6.8500000000000005,-.05){5};\draw[geoTrait,fill=otNavyTint](8.0,0)--++(.9,0)--++(-.3,.9)--++(-.3,-.35)--cycle;\node[below,font=\scriptsize]at(8.45,-.05){6};\draw[geoVecteur](0,1.25)--node[above]{$\vec u$ : 4 cm}(1.6,1.25);\end{tikzpicture}
{\par\smallskip\fontsize{8}{10}\selectfont\color{otMuted}Six motifs d’une frise ; chaque motif est l’image du précédent par la translation de vecteur $\vec u$.\par}
